\chapter{进程、同步与文件系统}

\section{进程是资源所有权集合}

进程包含地址空间、线程、打开文件、权限和生命周期状态。线程保存可调度执行
上下文。调度器在时钟中断或阻塞点选择可运行线程，并保存/恢复寄存器与栈。
抢占发生在任意允许中断的位置，因此共享状态必须有同步规则。

\section{锁、等待与 IPC}

自旋锁适合极短且不能睡眠的临界区；等待队列把不能继续的线程从运行队列移出；
管道把字节缓冲、读写端引用和唤醒条件组合成 IPC 对象。任何同步原语都要说明
谁能睡眠、谁能唤醒、锁顺序和中断上下文约束。

\section{从块到文件}

块设备只认识扇区。文件系统用超级块、inode、目录项和数据块提供名字与层次；
VFS 用统一对象接口连接进程文件描述符和具体文件系统。写操作还需要区分：
复制进缓存、提交给设备、设备报告完成、数据达到持久介质。

\section{失败与恢复}

磁盘可能在任意写入之间掉电。更新多个结构时，需要日志、写时复制或明确的
恢复算法保证旧状态或新状态至少一个完整。测试应注入越界块号、空间耗尽、
中途失败和损坏元数据，而不仅验证创建一个小文件。
